\documentclass[12pt,a4paper]{article}

\usepackage[utf8]{inputenc}
\usepackage[T1]{fontenc}
\usepackage[provide=*,spanish]{babel}
\usepackage{mathpazo}
\usepackage{microtype}
\usepackage{amsmath,amssymb,amsthm}
\usepackage{booktabs}
\usepackage{tabularx}
\usepackage{geometry}
\usepackage{enumitem}
\usepackage[hidelinks]{hyperref}

\geometry{margin=2.5cm}

\title{Ecuaciones Diferenciales Ordinarias\\Clase III -- Práctico}
\author{Benjamín Marcolongo}
\date{}

\begin{document}

	\maketitle

	\section*{Objetivo}

	En un problema de modelado no se comienza buscando una fórmula. Primero deben identificarse la magnitud desconocida, las leyes que describen su variación, los parámetros del modelo y las condiciones iniciales.

	En cada problema conviene seguir este esquema:

	\begin{enumerate}[label=\arabic*.]
		\item definir la variable dependiente y la variable independiente;
		\item elegir una convención de signos;
		\item traducir cada mecanismo a una tasa de cambio;
		\item comprobar la consistencia dimensional de la ecuación;
		\item incorporar las condiciones iniciales;
		\item analizar el dominio de validez del modelo.
	\end{enumerate}

	\clearpage
	\section{Modelos de crecimiento}

	Sea \(P(t)\) el tamaño de una población. Si durante un intervalo pequeño \(\Delta t\) el cambio es proporcional a la población presente,

	\[
	P(t+\Delta t)-P(t)
	\approx
	rP(t)\Delta t,
	\]
	%
	entonces, al dividir por \(\Delta t\) y tomar el límite \(\Delta t\to0\), obtenemos

	\[
	\boxed{\dot P=rP.}
	\]
	%
	El parámetro \(r\) tiene unidades de inversa de tiempo. Si \(r>0\), representa crecimiento neto; si \(r<0\), decrecimiento neto.

	\subsection{Inmigración constante}

	Si además ingresan \(I\) individuos por unidad de tiempo, la tasa total de cambio es

	\[
	\boxed{\dot P=rP+I.}
	\]
	%
	El término \(rP\) depende del estado del sistema, mientras que \(I\) representa un flujo externo constante. Las unidades deben satisfacer

	\[
	[r]=\frac{1}{\text{tiempo}},
	\qquad
	[I]=\frac{\text{individuos}}{\text{tiempo}}.
	\]

	\subsection{Natalidad lineal y mortalidad cuadrática}

	Supongamos que la tasa de natalidad es proporcional a \(P\), mientras que la tasa de mortalidad es proporcional a \(P^2\). Entonces

	\[
	\text{nacimientos por unidad de tiempo}=bP,
	\]
	%
	\[
	\text{muertes por unidad de tiempo}=dP^2,
	\]
	%
	y el balance resulta

	\[
	\boxed{\dot P=bP-dP^2=P(b-dP).}
	\]
	%
	El término cuadrático representa interacciones cuya frecuencia aumenta con la densidad poblacional. Si \(b,d>0\), el modelo puede escribirse como

	\[
	\dot P=bP\left(1-\frac{P}{K}\right),
	\qquad
	K=\frac{b}{d},
	\]
	%
	que es la forma logística con capacidad de carga \(K\).

	\clearpage
	\section{Ley de enfriamiento y calentamiento}

	Sea \(T(t)\) la temperatura de un cuerpo y \(T_m(t)\) la temperatura del medio. La ley de Newton establece que la tasa de cambio de \(T\) es proporcional a la diferencia entre la temperatura del cuerpo y la del ambiente:

	\[
	\boxed{\dot T=-k\bigl(T-T_m(t)\bigr),}
	\qquad k>0.
	\]
	%
	La elección del signo garantiza el comportamiento físico correcto:

	\begin{itemize}[leftmargin=*]
		\item si \(T>T_m\), entonces \(\dot T<0\) y el cuerpo se enfría;
		\item si \(T<T_m\), entonces \(\dot T>0\) y el cuerpo se calienta;
		\item si \(T=T_m\), entonces \(\dot T=0\) en ese instante.
	\end{itemize}

	Si la temperatura ambiente es constante, la ecuación es autónoma. Si \(T_m\) depende del tiempo, la ecuación deja de ser autónoma.

	Por ejemplo, una temperatura ambiente periódica con período de 24 horas puede modelarse mediante

	\[
	T_m(t)
	=
	10\sin\left(\frac{\pi}{12}t-\frac{\pi}{2}\right)+8.
	\]
	%
	En ese entorno, el modelo para el cuerpo es

	\[
	\boxed{
	\dot T
	=
	-k\left[
	T-10\sin\left(\frac{\pi}{12}t-\frac{\pi}{2}\right)-8
	\right].}
	\]
	%
	Para determinar una solución particular debe agregarse una condición inicial. Por ejemplo,

	\[
	T(t_0)=T_0.
	\]

	\clearpage
	\section{Drenaje de un tanque: ley de Torricelli}

	Sea \(h(t)\) la altura del líquido y \(A_w(h)\) el área de la sección transversal del tanque a esa altura. Un incremento \(dh\) produce un cambio de volumen

	\[
	dV=A_w(h)\,dh.
	\]
	%
	Por lo tanto,

	\[
	\frac{dV}{dt}=A_w(h)\dot h.
	\]
	%
	La ley de Torricelli aproxima la velocidad de salida por un orificio como

	\[
	v_s=\sqrt{2gh}.
	\]
	%
	Si el orificio tiene área \(a\) y se incorpora un coeficiente de descarga \(C_d\), el caudal saliente es

	\[
	Q_s=C_da\sqrt{2gh}.
	\]
	%
	Como el volumen disminuye,

	\[
	\frac{dV}{dt}=-Q_s.
	\]
	%
	El modelo general queda entonces

	\[
	\boxed{A_w(h)\dot h=-C_da\sqrt{2gh}.}
	\]

	\subsection{Tanque cónico}

	Consideremos un cono de altura total \(H\) y radio superior \(R\), con el vértice orientado hacia abajo. Por semejanza de triángulos,

	\[
	\frac{r(h)}{h}=\frac{R}{H},
	\]
	%
	de modo que

	\[
	r(h)=\frac{R}{H}h.
	\]

	El área de la sección transversal es

	\[
	A_w(h)
	=
	\pi r(h)^2
	=
	\pi\frac{R^2}{H^2}h^2.
	\]
	%
	Sustituyendo en el balance de volumen,

	\[
	\pi\frac{R^2}{H^2}h^2\dot h
	=
	-C_da\sqrt{2gh},
	\]
	%
	y, para \(h>0\),

	\[
	\boxed{
	\dot h
	=
	-\frac{C_daH^2\sqrt{2g}}{\pi R^2}\,h^{-3/2}.}
	\]
	%
	\enlargethispage{3\baselineskip}
	La geometría del tanque modifica la relación entre la disminución del volumen y la disminución de la altura. Por eso no puede suponerse un área transversal constante.

	\clearpage
	\section{Movimiento con masa variable}

	Consideremos un cable uniforme de densidad lineal \(\rho\). Sea \(x(t)\) la longitud de cable que se encuentra en movimiento y tomemos como positiva la dirección vertical hacia abajo.

	La masa del tramo móvil es

	\[
	m(t)=\rho x(t),
	\]
	%
	y su velocidad es

	\[
	v(t)=\dot x(t).
	\]

	Siguiendo el balance de cantidad de movimiento propuesto para el sistema,

	\[
	F=\frac{d}{dt}\bigl(mv\bigr).
	\]
	%
	La fuerza externa que impulsa el tramo móvil es su peso,

	\[
	F=m(t)g=\rho xg.
	\]
	%
	Como

	\[
	mv=\rho x\dot x,
	\]
	%
	la regla del producto da

	\[
	\frac{d}{dt}\bigl(\rho x\dot x\bigr)
	=
	\rho\left(\dot x^2+x\ddot x\right).
	\]
	%
	Por lo tanto,

	\[
	\rho\left(\dot x^2+x\ddot x\right)
	=
	\rho xg,
	\]
	%
	y el modelo se reduce a

	\[
	\boxed{x\ddot x+\dot x^2=gx.}
	\]
	%
	Si inicialmente cuelga una longitud \(x_0>0\) y el cable está en reposo, las condiciones iniciales son

	\[
	\boxed{x(0)=x_0,
	\qquad
	\dot x(0)=0.}
	\]

	La densidad \(\rho\) se cancela porque tanto la fuerza gravitatoria como la cantidad de movimiento son proporcionales a ella. La ecuación obtenida describe el modelo bajo las hipótesis indicadas; en un sistema de masa variable, la elección de la frontera del sistema y el flujo de cantidad de movimiento deben especificarse cuidadosamente.

	\clearpage
	\section*{Articulación con la Guía 1}

	Los modelos desarrollados corresponden a los problemas 10--13 de la Guía 1. En esta etapa el objetivo principal es construir correctamente las ecuaciones diferenciales. Su resolución se realizará con los métodos analíticos que se incorporen al curso.

	\subsection*{Nota sobre la elaboración del material}

	Este material fue elaborado por Benjamín Marcolongo con la asistencia de \textbf{GPT-5.6 Sol}, un modelo de OpenAI utilizado a través del entorno Codex, para la organización, redacción y revisión del contenido.

\end{document}
